% ═══════════════════════════════════════════════════════════════════════════
% QR-CODES: Ohmsches Gesetz (Physik Klasse 8, Doppelstunde 02a)
% Enthält: Lernpfad + 4 Simulationen
% ═══════════════════════════════════════════════════════════════════════════

\documentclass[a4paper,11pt]{article}

\usepackage[utf8]{inputenc}
\usepackage[T1]{fontenc}
\usepackage[ngerman]{babel}
\usepackage{helvet}
\renewcommand{\familydefault}{\sfdefault}

\usepackage[margin=15mm]{geometry}
\usepackage{tikz}
\usepackage{xcolor}
\usepackage{qrcode}

% ═══════════════════════════════════════════════════════════════════════════
% FARBEN
% ═══════════════════════════════════════════════════════════════════════════

\definecolor{physik}{HTML}{2c5aa0}
\colorlet{fachfarbe}{physik}

% ═══════════════════════════════════════════════════════════════════════════
% URLs
% ═══════════════════════════════════════════════════════════════════════════

\newcommand{\baseurl}{https://devbox42.github.io/sciencesim/physik/kl8-ohm}

% ═══════════════════════════════════════════════════════════════════════════
% DOKUMENT
% ═══════════════════════════════════════════════════════════════════════════

\pagestyle{empty}

\begin{document}

% ───────────────────────────────────────────────────────────────────────────
% SEITE 1: Lernpfad (groß, zum Aufhängen)
% ───────────────────────────────────────────────────────────────────────────

\begin{center}

\vspace*{5mm}
{\fontsize{26}{32}\selectfont\bfseries\color{fachfarbe}Lernpfad: Das Ohmsche Gesetz}

\vspace{12mm}

\fcolorbox{fachfarbe}{white}{\qrcode[height=95mm]{\baseurl/LP-02a-ohmsches-gesetz.html}}

\vspace{10mm}

{\large\color{gray}devbox42.github.io/.../LP-02a-ohmsches-gesetz.html}

\vspace{8mm}

{\Large\color{gray}Scanne den QR-Code mit deinem Gerät}

\vspace{8mm}

\fcolorbox{fachfarbe}{fachfarbe!10}{%
    \parbox{0.75\textwidth}{%
        \centering
        \vspace{3mm}
        Arbeite die 16 Schritte durch und fülle dabei\\[2mm]
        das Arbeitsblatt AB-02a aus.
        \vspace{3mm}
    }%
}

\vspace{10mm}
{\footnotesize\color{gray}Physik Klasse 8 | Doppelstunde 02a | 80 Min}

\end{center}

\newpage

% ───────────────────────────────────────────────────────────────────────────
% SEITE 2: Alle 4 Simulationen (zum Ausdrucken für Tische)
% ───────────────────────────────────────────────────────────────────────────

\begin{center}

\vspace*{3mm}
{\fontsize{20}{24}\selectfont\bfseries\color{fachfarbe}Simulationen: Ohmsches Gesetz}

\vspace{6mm}

% ─── Simulation 1: Wassermodell ───
\fcolorbox{fachfarbe}{white}{%
    \begin{minipage}{0.22\textwidth}
        \centering
        \vspace{2mm}
        {\scriptsize\bfseries\color{fachfarbe}Wassermodell}\\[3mm]
        \qrcode[height=32mm]{\baseurl/SIM-02a-wassermodell.html}\\[2mm]
        {\tiny\color{gray}SIM-02a-wassermodell}
        \vspace{2mm}
    \end{minipage}%
}
\hfill
% ─── Simulation 2: Kennlinie ───
\fcolorbox{fachfarbe}{white}{%
    \begin{minipage}{0.22\textwidth}
        \centering
        \vspace{2mm}
        {\scriptsize\bfseries\color{fachfarbe}I(U)-Kennlinie}\\[3mm]
        \qrcode[height=32mm]{\baseurl/SIM-02a-kennlinie.html}\\[2mm]
        {\tiny\color{gray}SIM-02a-kennlinie}
        \vspace{2mm}
    \end{minipage}%
}
\hfill
% ─── Simulation 3: Glühlampe ───
\fcolorbox{fachfarbe}{white}{%
    \begin{minipage}{0.22\textwidth}
        \centering
        \vspace{2mm}
        {\scriptsize\bfseries\color{fachfarbe}Glühlampe}\\[3mm]
        \qrcode[height=32mm]{\baseurl/SIM-02a-gluehlampe.html}\\[2mm]
        {\tiny\color{gray}SIM-02a-gluehlampe}
        \vspace{2mm}
    \end{minipage}%
}
\hfill
% ─── Simulation 4: Atomrümpfe ───
\fcolorbox{fachfarbe}{white}{%
    \begin{minipage}{0.22\textwidth}
        \centering
        \vspace{2mm}
        {\scriptsize\bfseries\color{fachfarbe}Atomrümpfe}\\[3mm]
        \qrcode[height=32mm]{\baseurl/SIM-02a-atomruempfe.html}\\[2mm]
        {\tiny\color{gray}SIM-02a-atomruempfe}
        \vspace{2mm}
    \end{minipage}%
}

\vspace{10mm}

% ─── Schneidemarken ───
{\color{gray}\hrule}

\vspace{6mm}

% ─── Einzelne Kärtchen zum Ausschneiden (4×) ───
{\small\color{fachfarbe}\bfseries Zum Ausschneiden – einzelne QR-Kärtchen}

\vspace{6mm}

% Zeile 1: Wassermodell + Kennlinie
\fcolorbox{fachfarbe}{white}{%
    \begin{minipage}{45mm}
        \centering
        \vspace{2mm}
        {\footnotesize\bfseries\color{fachfarbe}Wassermodell}\\[2mm]
        \qrcode[height=28mm]{\baseurl/SIM-02a-wassermodell.html}\\[1mm]
        {\tiny Physik Kl. 8 | Ohm}
        \vspace{2mm}
    \end{minipage}%
}
\hspace{3mm}
\fcolorbox{fachfarbe}{white}{%
    \begin{minipage}{45mm}
        \centering
        \vspace{2mm}
        {\footnotesize\bfseries\color{fachfarbe}I(U)-Kennlinie}\\[2mm]
        \qrcode[height=28mm]{\baseurl/SIM-02a-kennlinie.html}\\[1mm]
        {\tiny Physik Kl. 8 | Ohm}
        \vspace{2mm}
    \end{minipage}%
}
\hspace{3mm}
\fcolorbox{fachfarbe}{white}{%
    \begin{minipage}{45mm}
        \centering
        \vspace{2mm}
        {\footnotesize\bfseries\color{fachfarbe}Glühlampe}\\[2mm]
        \qrcode[height=28mm]{\baseurl/SIM-02a-gluehlampe.html}\\[1mm]
        {\tiny Physik Kl. 8 | Ohm}
        \vspace{2mm}
    \end{minipage}%
}
\hspace{3mm}
\fcolorbox{fachfarbe}{white}{%
    \begin{minipage}{45mm}
        \centering
        \vspace{2mm}
        {\footnotesize\bfseries\color{fachfarbe}Atomrümpfe}\\[2mm]
        \qrcode[height=28mm]{\baseurl/SIM-02a-atomruempfe.html}\\[1mm]
        {\tiny Physik Kl. 8 | Ohm}
        \vspace{2mm}
    \end{minipage}%
}

\vspace{10mm}

{\color{gray}\hrule}

\vspace{6mm}

{\footnotesize\color{gray}Seite 1: groß aufhängen | Seite 2: für Gruppentische ausschneiden}

\end{center}

\end{document}
